\documentclass[12pt]{article}

\usepackage{algorithm}
\usepackage{algpseudocode}
\usepackage{amsthm}
\usepackage{amsmath}
\usepackage{amsfonts}
\usepackage{mathrsfs}
\usepackage{amssymb}
\usepackage{units}
\usepackage{graphicx}
\usepackage{tikz-cd}
\usepackage{nicefrac}
\usepackage[hidelinks]{hyperref}
\usepackage{bbm}
\usepackage{color}
\usepackage{tensor}
\usepackage{tipa}
\usepackage{bussproofs}
\usepackage{ stmaryrd }
\usepackage{ textcomp }
\usepackage{leftidx}
\usepackage{afterpage}
\usepackage{varwidth}

\newcommand\blankpage{
    \null
    \thispagestyle{empty}
    \addtocounter{page}{-1}
    \newpage
    }

\graphicspath{ {images/} }

\theoremstyle{plain}
\newtheorem{thm}{Theorem}[subsection] % reset theorem numbering for each chapter
\newtheorem{proposition}[thm]{Proposition}
\newtheorem{lemma}[thm]{Lemma}
\newtheorem{fact}[thm]{Fact}
\newtheorem{cor}[thm]{Corollary}
\newtheorem{claim}[thm]{Claim}

\theoremstyle{definition}
\newtheorem{defn}[thm]{Definition} % definition numbers are dependent on theorem numbers
\newtheorem{exmp}[thm]{Example} % same for example numbers
\newtheorem{notation}[thm]{Notation}
\newtheorem{remark}[thm]{Remark}
\newtheorem{condition}[thm]{Condition}
\newtheorem{question}[thm]{Question}
\newtheorem{construction}[thm]{Construction}
\newtheorem{exercise}[thm]{Exercise}
\newtheorem{example}[thm]{Example}
\newtheorem{aside}[thm]{Aside}
\newtheorem{solution}[thm]{Solution}
%\newtheorem{algorithm}[thm]{Algorithm}

\newcommand{\bb}[1]{\mathbb{#1}}
\newcommand{\scr}[1]{\mathscr{#1}}
\newcommand{\call}[1]{\mathcal{#1}}
\newcommand{\psheaf}{\text{\underline{Set}}^{\scr{C}^{\text{op}}}}
\newcommand{\und}[1]{\underline{\hspace{#1 cm}}}
\newcommand{\adj}[1]{\text{\textopencorner}{#1}\text{\textcorner}}
\newcommand{\comment}[1]{}
\newcommand{\lto}{\longrightarrow}
\newcommand{\Spec}{\operatorname{Spec}}
\newcommand{\Proj}{\operatorname{Proj}}
\newcommand{\im}{\operatorname{im}}
\newcommand{\height}{\operatorname{ht}}
\DeclareMathOperator{\Hom}{Hom}

\title{Introduction to Varieties and Hartshorne Exercise Solutions}
\author{Will Troiani}
\date{Corrected and combined version}

\begin{document}

\maketitle
\tableofcontents
\section{Introduction}
We need the Hilbert scheme to model the exponential connective in linear logic. These notes begin the trek there.

\subsection*{Conventions}
Throughout, $k$ denotes an algebraically closed field unless explicitly stated otherwise. An \emph{algebraic set} means an arbitrary Zariski-closed subset of affine or projective space. An affine or projective \emph{variety} means an irreducible algebraic set with the induced topology. A quasi-affine, respectively quasi-projective, variety is a nonempty open subset of an affine, respectively projective, variety. With these conventions, coordinate rings of affine varieties are integral domains, and $K(X)$ denotes the function field of an irreducible variety $X$.
\section{Affine Varieties}
\subsection{Algebraic sets and the ideal of a set}
Throughout, we work with an algebraically closed field $k$.
\begin{defn}
Let $\mathfrak{a} \subseteq k[x_1,...,x_n]$ be an ideal. Define the \textbf{zero set of $\mathfrak{a}$}:
\begin{equation}
    Z(\mathfrak{a}) := \lbrace (p_1,...,p_n) \in k^n \mid \forall q \in \mathfrak{a}, q(p_1,...,p_n) = 0\rbrace
\end{equation}
A subset $T \subseteq k^n$ for which there exists an ideal $\mathfrak{a} \subseteq k[x_1,...,x_n]$ so that $T = Z(\mathfrak{a})$ is an \textbf{algebraic set}.

Given an arbitrary set of points $T \subseteq k^n$ we define the \textbf{ideal of $T$}:
\begin{equation}
    I(T) := \lbrace q \in k[x_1,...,x_n] \mid \forall p = (p_1,...,p_n) \in T, q(p_1,...,p_n) = 0 \rbrace
\end{equation}
which is in fact an ideal.
\end{defn}
Two non-equal ideals may have equal zero sets, for instance $Z(x^2) = Z(x)$ (recall $k$ is a field).
\subsubsection{When do two ideals yield the same zero set?}
This will be answered by \emph{Hilbert's Nullstellensatz}, which we state on the level of generality where $k$ is not necessarily algebraically closed.
\begin{defn}
\label{def:algzero}
Let $F$ be a field (not necessarily algebraically closed). An \textbf{algebraic zero} of a subset $\Phi \subseteq F[x_1,...,x_n]$ is a sequence $(\alpha_1,...,\alpha_n)$ of elements in an algebraic closure $\bar{F}$ such that $f(\alpha_1,...,\alpha_n) = 0$ for all $f \in \Phi$.
\end{defn}
Notice that if a root exists in any algebraic closure it exists in them all, so it makes sense to talk about an algebraic zero in absence of a particular algebraic closure.

Given a field $F$ (not necessarily algebraically closed) a set $\Phi \subseteq F[x_1, \ldots, x_n]$ we write $(\Phi)$ for the ideal generated by the elements of $\Phi$.
\begin{thm}
\label{thm:hilbertsnullstellensatz}
Let $\Phi \subseteq F[x_1,...,x_n]$.
\begin{enumerate}
    \item\label{generates} If $\Phi$ admits no algebraic zeros, then $(\Phi) = F[x_1,...,x_n]$.
    \item\label{power} Let $f \in F[x_1,...,x_n]$ be such that $f(\alpha_1,...,\alpha_n) = 0$ for all algebraic zeros $(\alpha_1,...,\alpha_n)$ of $\Phi$, then there exists $r > 0$ such that $f^r \in (\Phi)$.
\end{enumerate}
\end{thm}
\begin{proof}
See \cite{algebra}.
\end{proof}
Thus we have:
\begin{lemma}
Let $\mathfrak{a} \subseteq k[x_1,...,x_n]$ be an ideal. Then
\begin{equation}\label{eq:IZ}
    IZ(\mathfrak{a}) = \sqrt{\mathfrak{a}}
\end{equation}
\end{lemma}
We thus have $ZIZ(\mathfrak{a}) = Z(\sqrt{\mathfrak{a}})$.  It remains to show $Z(\sqrt{\mathfrak{a}}) = Z(\mathfrak{a})$, this will require a new perspective on $Z(\mathfrak{a})$:
\begin{lemma}
The algebraic sets form a collection of closed sets, in fact:
\begin{enumerate}
    \item $\varnothing = Z(1)$, and $k^n = Z(0)$,
    \item $Z(\mathfrak{a}) \cup Z(\mathfrak{b}) = Z(\mathfrak{a}\mathfrak{b})$,
    \item $\bigcap_{i \in I}Z(\mathfrak{a}_i) = Z(\mathfrak{b})$, where $\mathfrak{b}$ is the ideal generated by $\bigcup_{i \in I}\mathfrak{a}_i$.
\end{enumerate}
\end{lemma}
From now on, we consider $k^n$ as a topological space whose closed sets are given by the algebraic sets, this is the \textbf{Zariski Topology}.
\begin{lemma}\label{lem:closure_zt}
Let $T\subseteq \bb A^n$. Then
\[
ZI(T)=\overline T.
\]
\end{lemma}
\begin{proof}
We have $T\subseteq ZI(T)$ and $ZI(T)$ is closed. Conversely, suppose $T\subseteq W=Z(\mathfrak a)$ with $W$ closed. Then $\mathfrak a\subseteq I(T)$, hence $ZI(T)\subseteq Z(\mathfrak a)=W$. Thus $ZI(T)$ is the smallest closed set containing $T$, i.e. $ZI(T)=\overline T$.
\end{proof}
We thus have $ZIZ(\mathfrak{a}) = \overline{Z(\mathfrak{a})} = Z(\mathfrak{a})$.
%
%
%

\subsection{Irreducible sets}\label{sec:irreducible}
In Section \ref{sec:dimension} the dimension of a vector space will be generalised to make sense for arbitrary affine varieties.
\begin{lemma}\label{lem:irred_criterion}
Let $Y\subseteq X$ be a subset of a topological space. Then the following are equivalent:
\begin{enumerate}
\item $Y$ is irreducible;
\item every non-empty open subset of $Y$ is dense in $Y$;
\item every pair of non-empty open subsets of $Y$ has non-empty intersection.
\end{enumerate}
\end{lemma}
\begin{proof}
If $Y$ is irreducible and $U$ is a non-empty open subset which is not dense, then $Y=\overline U\cup(Y\setminus U)$ is a union of two proper closed subsets. The implication from density to pairwise intersection is immediate. Conversely, if $Y=F_1\cup F_2$ with $F_i$ proper closed, then $Y\setminus F_1$ and $Y\setminus F_2$ are disjoint non-empty open subsets.
\end{proof}
\begin{lemma}\label{lem:irred_subspace_invariance}
If $Z\subseteq Y\subseteq X$, then $Z$ is irreducible with the subspace topology inherited from $X$ if and only if it is irreducible with the subspace topology inherited from $Y$.
\end{lemma}
\begin{proof}
The two subspace topologies on $Z$ are identical.
\end{proof}
\begin{lemma}\label{lem:open_irred_dense}
Every non-empty open subset of an irreducible space is irreducible and dense.
\end{lemma}
\begin{proof}
Density is Lemma \ref{lem:irred_criterion}. If $U=F_1\cup F_2$ with $F_i$ closed in $U$, write $F_i=U\cap E_i$ with $E_i$ closed in the ambient irreducible space. The density of $U$ forces the ambient space to be $E_1\cup E_2$, so one $E_i$ is the whole space and one $F_i$ is all of $U$.
\end{proof}
\begin{lemma}\label{lem:closure_irred}
The closure of an irreducible set is irreducible.
\end{lemma}
\begin{proof}
If $\overline Y=F\cup G$ with $F,G$ closed in $\overline Y$, then $Y=(Y\cap F)\cup(Y\cap G)$. Irreducibility of $Y$ puts $Y$ in one of $F,G$; taking closures puts $\overline Y$ there too.
\end{proof}

\subsection{Dimension}\label{sec:dimension}
The dimension of a variety is defined by chains of irreducible closed subsets. For affine space there is the chain
\begin{equation}\label{eq:dimension_chain}
\bb A^n=Z(0)\supsetneq Z(x_n)\supsetneq\cdots\supsetneq Z(x_2,\ldots,x_n)\supsetneq Z(x_1,\ldots,x_n),
\end{equation}
which has length $n$.
\begin{defn}
An affine variety $L\subseteq\bb A^n$ is \textbf{affine linear} if it is cut out by linear polynomials. Equivalently, $L$ is a translate of a vector subspace of $k^n$.
\end{defn}
An affine linear subspace of dimension $m$ is isomorphic to $\bb A^m$, so it has dimension $m$.
\begin{defn}\label{def:dimension_var}
The \textbf{dimension} of a topological space $X$ is the supremum of the lengths $n$ of strictly descending chains
\[
X_0\supsetneq X_1\supsetneq\cdots\supsetneq X_n
\]
of irreducible closed subsets. The dimension of a quasi-affine variety is its dimension as a topological space.
\end{defn}
\begin{lemma}\label{lem:the_lemma_never_stated}
Let $X$ be a topological space and $U\subseteq X$ open. There is a strict-inclusion-preserving bijection
\[
\{\text{irreducible closed }C\subseteq X:C\cap U\neq\varnothing\}
\longleftrightarrow
\{\text{irreducible closed }D\subseteq U\},
\]
given by $C\mapsto C\cap U$ and $D\mapsto\overline D^{X}$.
\end{lemma}
\begin{proof}
If $C$ is irreducible and meets $U$, then $C\cap U$ is a non-empty open subset of $C$, hence irreducible and dense in $C$. If $D$ is irreducible closed in $U$, then $\overline D^X$ is irreducible and $\overline D^X\cap U=D$. These assignments are inverse and preserve strict inclusions.
\end{proof}
\begin{cor}\label{cor:dimension_open_cover}
If $X$ is covered by open subsets $U_i$, then $\dim X=\sup_i\dim U_i$.
\end{cor}
\begin{proof}
Every chain in $U_i$ closes to a chain in $X$. Conversely, a chain in $X$ meets some $U_i$ at its smallest member, and Lemma \ref{lem:the_lemma_never_stated} restricts it to a chain of the same length in $U_i$.
\end{proof}
\begin{thm}\label{thm:onepointeight}
Let $A$ be a finitely generated $k$-domain. Then
\begin{enumerate}
\item $\operatorname{tr.deg}_k A=\dim A$;
\item if $\mathfrak p$ is prime in $A$, then $\operatorname{ht}\mathfrak p+\dim A/\mathfrak p=\dim A$.
\end{enumerate}
\end{thm}
\begin{proof}
See \cite{algebra}.
\end{proof}
\begin{cor}\label{lem:quasi_dimension}
Let $Y$ be a quasi-affine variety. Then $\dim Y=\dim\overline Y$.
\end{cor}
\begin{proof}
A quasi-affine variety is a non-empty open subset of its irreducible affine closure. Non-empty open subsets of irreducible varieties have the same dimension by Lemma \ref{lem:the_lemma_never_stated} and the chain argument above.
\end{proof}
\section{Projective varieties}
%
\begin{defn}
Consider the set $\bb{A}^{n+1}\setminus\lbrace 0 \rbrace$ modded out by the equivalence relation $x \sim \lambda x$ for any $\lambda \in k\setminus \lbrace 0 \rbrace$. We denote this set $\bb{P}^n$. This is endowed with a topological space but \emph{not} the quotient space topology, instead the closed sets are defined to be such zero sets of homogeneous polynomials (polynomials which are the sum of monomials of the same degree). \textbf{Projective $n$-space} is $\bb{P}^n$ along with this topology. (See \cite{hartshorne} for more details).
\end{defn}
\begin{defn}
A \textbf{Projective variety} is an irreducible, closed subset of $\bb{P}^n$. A \textbf{quasi-projective variety} is an open subset of a projective variety.
\end{defn}
\begin{notation}\label{notation:affine_proj}
If extra clarity is needed, given a homogeneous ideal $\mathfrak{a} \subseteq k[x_0,...,x_n]$, the notation $Z_{\bb{P}^n}(\mathfrak{a})$ will be used for the zero set of $\mathfrak{a}$ in $\bb{P}^n$ and $Z_{\bb{A}^{n+1}}(\mathfrak{a})$ for the zero set $Z(\hat{\mathfrak{a}})$ in $\bb{A}^{n+1}$ where $\hat{\mathfrak{a}} \subseteq \bb{A}^{n+1}$ is defined by
\begin{equation*}
    f(x_1,...,x_{n+1}) \in \hat{\mathfrak{a}} \subseteq k[x_1,...,x_{n+1}] \Longleftrightarrow f(x_0,...,x_n) \in \mathfrak{a} \subseteq k[x_0,...,x_n]
\end{equation*}
\end{notation}
In practice, projective varieties are dealt with by considering a related affine variety, one way of doing this is the following: given a a non-empty projective variety $\varnothing \neq Z_{\bb{P}^n}(\mathfrak{a}) \subseteq \bb{P}^n$, consider $Z_{\bb{A}^{n+1}}(\mathfrak{a})\subseteq \bb{A}^{n+1}$ (as per Notation \ref{notation:affine_proj}), which consists of all representations of equivalence classes in $Z(\mathfrak{a}) \subseteq \bb{P}^n$ along with, by homogeneity of $\mathfrak{a}$, a new point $0$. One then deals with these two situations separately. See the solution to Exercise $2.1,2.2,2.3$ \cite{hartshorne_solutions} for two examples of this.

Another way of relating a projective variety to a related affine variety is by using the fact that projective varieties are covered homeomorphically by affine varieties, a result the next Section \ref{sec:homog_proj} is dedicated to proving.
\subsection{Homogenisation and projection}\label{sec:homog_proj}
Projective space is covered by affine charts. Let
\[
U_l=\bb P^n\setminus Z(x_l).
\]
Then
\begin{align*}
\varphi_l:U_l&\longrightarrow \bb A^n,\\
[a_0:\cdots:a_n]&\longmapsto (a_0/a_l,\ldots,\widehat{a_l/a_l},\ldots,a_n/a_l)
\end{align*}
is a homeomorphism, with inverse obtained by inserting a $1$ in the $l$-th coordinate. The proof is the geometric form of dehomogenisation and homogenisation.

\begin{defn}\label{def:projection}
Let $F\in k[x_0,\ldots,x_n]$ be homogeneous. The \textbf{$l$-th dehomogenisation} of $F$ is
\[
F_{(l)}:=F(x_0,\ldots,x_{l-1},1,x_{l+1},\ldots,x_n),
\]
viewed as a polynomial in the $n$ remaining variables. If $\mathfrak b\subseteq k[x_0,\ldots,x_n]$ is homogeneous, let $\mathfrak b_{(l)}$ be the ideal generated by all $F_{(l)}$ with $F\in\mathfrak b$ homogeneous.
\end{defn}

\begin{defn}\label{def:homogenisation}
Let $f\in k[x_1,\ldots,x_n]$ have degree at most $d$. Its \textbf{homogenisation to degree $d$ with respect to $x_0$} is
\[
f^{h,d}(x_0,x_1,\ldots,x_n)=x_0^d f(x_1/x_0,\ldots,x_n/x_0).
\]
When $d=\deg f$, write $f^h$. For an ideal $\mathfrak a\subseteq k[x_1,\ldots,x_n]$, let $\mathfrak a^h$ denote the homogeneous ideal generated by all $f^h$ with $f\in\mathfrak a$.
\end{defn}

\begin{lemma}\label{lem:poly_homog_bij}
For each $d\geq0$, dehomogenisation at $x_0$ gives a bijection between homogeneous polynomials of degree $d$ in $k[x_0,\ldots,x_n]$ and polynomials in $k[x_1,\ldots,x_n]$ of degree at most $d$.
\end{lemma}
\begin{proof}
If $F$ is homogeneous of degree $d$, then
\[
F=x_0^dF(1,x_1/x_0,\ldots,x_n/x_0),
\]
so $F$ is recovered from its dehomogenisation by homogenising back to degree $d$. The other composite is immediate after setting $x_0=1$.
\end{proof}

\begin{lemma}\label{closedsets}
For homogeneous ideals $\mathfrak b\subseteq k[x_0,\ldots,x_n]$ and ideals $\mathfrak a\subseteq k[x_1,\ldots,x_n]$,
\begin{align*}
\varphi_l\big(U_l\cap Z_{\bb P^n}(\mathfrak b)\big)&=Z_{\bb A^n}(\mathfrak b_{(l)}),\\
\varphi_l^{-1}\big(Z_{\bb A^n}(\mathfrak a)\big)&=U_l\cap Z_{\bb P^n}(\mathfrak a^h).
\end{align*}
\end{lemma}
\begin{proof}
Both identities are obtained by setting $x_l=1$ on the chart $U_l$ and homogenising back to the relevant degree. The fixed-degree condition in Lemma \ref{lem:poly_homog_bij} is essential; without fixing the degree, dehomogenisation is not injective.
\end{proof}

\subsection{Gr\"{o}bner Bases}
Given generators $f_1,\ldots,f_m\in k[x_1,\ldots,x_n]$ for an ideal $I$, homogenising the displayed generators need not generate the whole homogenised ideal. For instance, if $f_1=x^2-y$ and $f_2=x^3-z$, then $-xf_1+f_2=z-xy$, so the homogenisation $x_0z-xy$ belongs to the homogenised ideal even though it is not generated by $x^2-x_0y$ and $x^3-x_0^2z$. Gr\"obner bases provide a systematic way to choose generators which behave well under leading terms and, for degree-compatible orders, under homogenisation.

\begin{defn}
A \textbf{monomial order} on $k[x_1,\ldots,x_n]$ is a total well-ordering of the monomials such that $m_1<m_2$ implies $m_1m<m_2m$ for every monomial $m$. For a nonzero polynomial $f$, write $\operatorname{LT}(f)$ for its leading term and $\operatorname{LM}(f)$ for its leading monomial.
\end{defn}

\begin{defn}
Let $I\subseteq k[x_1,\ldots,x_n]$ be an ideal. A finite set $G=\{g_1,\ldots,g_r\}\subseteq I$ is a \textbf{Gr\"obner basis} of $I$ if
\[
(\operatorname{LT}(I))=(\operatorname{LT}(g_1),\ldots,\operatorname{LT}(g_r)).
\]
Such a set automatically generates $I$.
\end{defn}

\begin{lemma}[Dickson's Lemma]\label{lem:Dickson}
Every monomial ideal in $k[x_1,\ldots,x_n]$ is generated by finitely many monomials.
\end{lemma}
\begin{proof}
The standard induction proof reduces from $n$ variables to $n-1$ variables. For a monomial ideal $I\subseteq k[x_1,\ldots,x_n]$, let $J_m\subseteq k[x_1,\ldots,x_{n-1}]$ be the monomial ideal consisting of monomials $u$ such that $ux_n^m\in I$. Then
\[
J_0\subseteq J_1\subseteq J_2\subseteq\cdots.
\]
By the induction hypothesis and Noetherianity in $n-1$ variables, this ascending chain stabilises, say at $J_N$. Choose finite monomial generators for $J_0,\ldots,J_N$. The finitely many monomials $ux_n^m$, with $u$ among the chosen generators of $J_m$ and $0\leq m\leq N$, generate $I$.
\end{proof}

\begin{cor}\label{cor:grobner_basis_existence}
Every ideal $I\subseteq k[x_1,\ldots,x_n]$ admits a finite Gr\"obner basis.
\end{cor}
\begin{proof}
The ideal $(\operatorname{LT}(I))$ is a monomial ideal, hence by Dickson's Lemma it is generated by finitely many leading terms $\operatorname{LT}(g_1),\ldots,\operatorname{LT}(g_r)$ with $g_i\in I$. Then $g_1,\ldots,g_r$ is a Gr\"obner basis of $I$.
\end{proof}

\begin{defn}
Let $f,g\in k[x_1,\ldots,x_n]$ be nonzero. If $\operatorname{LM}(f)=x^\alpha$ and $\operatorname{LM}(g)=x^\beta$, set $\gamma_i=\max(\alpha_i,\beta_i)$, so that $x^\gamma=\operatorname{lcm}(\operatorname{LM}(f),\operatorname{LM}(g))$. The \textbf{$S$-polynomial} of $f$ and $g$ is
\[
S(f,g)=\frac{x^\gamma}{\operatorname{LT}(f)}f-\frac{x^\gamma}{\operatorname{LT}(g)}g.
\]
\end{defn}

\begin{lemma}[Buchberger criterion]\label{lem:suffices_S_poly}
A generating set $g_1,\ldots,g_r$ of an ideal $I$ is a Gr\"obner basis if and only if every $S(g_i,g_j)$ reduces to zero on division by $g_1,\ldots,g_r$.
\end{lemma}
\begin{proof}
See \cite[\S 2.6, Theorem 6]{grobner}.
\end{proof}

\begin{algorithm}\label{alg:division_alg}
\caption{Multivariable division algorithm}\label{alg:division}
\begin{algorithmic}
\Require $(f_1,\ldots,f_s),f$
\State $q_1,\ldots,q_s\gets 0,\ldots,0$
\State $r\gets 0$
\State $p\gets f$
\While{$p\neq 0$}
\State $\text{DivOcc}\gets\texttt{False}$
\State $i\gets 1$
\While{$i\leq s$ and $\text{DivOcc}=\texttt{False}$}
\If{$\operatorname{LT}(f_i)$ divides $\operatorname{LT}(p)$}
\State $a\gets \operatorname{LT}(p)/\operatorname{LT}(f_i)$
\State $q_i\gets q_i+a$
\State $p\gets p-af_i$
\State $\text{DivOcc}\gets\texttt{True}$
\Else
\State $i\gets i+1$
\EndIf
\EndWhile
\If{$\text{DivOcc}=\texttt{False}$}
\State $r\gets r+\operatorname{LT}(p)$
\State $p\gets p-\operatorname{LT}(p)$
\EndIf
\EndWhile
\State \Return{$(q_1,\ldots,q_s,r)$}
\end{algorithmic}
\end{algorithm}

For the homogenisation problem above we use the following standard theorem instead of reproving it here.
\begin{thm}\label{thm:homogenized_grobner}
Let $G$ be a Gr\"obner basis of an ideal $I\subseteq k[x_1,\ldots,x_n]$ for a degree-compatible monomial order. With the corresponding homogenised order on $k[x_0,x_1,\ldots,x_n]$, the set $G^h=\{g^h:g\in G\}$ is a Gr\"obner basis of the homogenised ideal $I^h$.
\end{thm}
\begin{proof}
This is the standard homogenisation theorem for Gr\"obner bases; see \cite[Chapter 2]{grobner}.
\end{proof}
As an application, to generate the homogenised ideal $I^h$, first extend any generating set of $I$ to a Gr\"obner basis for a degree-compatible order, and then homogenise that Gr\"obner basis.

\subsection{Segre embedding}\label{subsec:segre}
The \textbf{Segre embedding} is
\[
\psi:\bb P^r\times\bb P^s\longrightarrow \bb P^{(r+1)(s+1)-1},\qquad
([P_i],[Q_j])\longmapsto [P_iQ_j]_{0\leq i\leq r,\,0\leq j\leq s}.
\]
Here $\times$ denotes the product variety, not the product topology.  Let
\[
\theta:k[z_{ij}:0\leq i\leq r,0\leq j\leq s]\longrightarrow k[x_0,\ldots,x_r,y_0,\ldots,y_s],\qquad z_{ij}\longmapsto x_i y_j.
\]
The kernel of $\theta$ is generated by the $2\times2$ minors
\[
z_{ij}z_{kl}-z_{il}z_{kj}.
\]
\begin{lemma}\label{lem:segre_image}
The image of the Segre embedding is $Z(\ker\theta)$, and the Segre map is injective.
\end{lemma}
\begin{proof}
The inclusion $\im\psi\subseteq Z(\ker\theta)$ is immediate. Conversely, let $P=[P_{ij}]\in Z(\ker\theta)$ and choose $(a,b)$ with $P_{ab}\neq0$. The $2\times2$ minors give
\[
P_{ij}P_{ab}=P_{ib}P_{aj}
\]
for all $i,j$, hence
\[
P_{ij}=\frac{P_{ib}P_{aj}}{P_{ab}}.
\]
Thus $P$ is the Segre image of
\[
[P_{0b}:\cdots:P_{rb}],\qquad [P_{a0}:\cdots:P_{as}].
\]
This also proves injectivity.
\end{proof}
\begin{proposition}\label{prop:segre_isomorphism}
The Segre embedding is an isomorphism onto its closed image.
\end{proposition}
\begin{proof}
On the affine chart $z_{ab}\neq0$ of the image, the inverse is given by the regular formula
\[
[P_{ij}]\longmapsto ([P_{0b}:\cdots:P_{rb}], [P_{a0}:\cdots:P_{as}]).
\]
These formulae cover the image and agree on overlaps. Thus the Segre map has a regular inverse onto $Z(\ker\theta)$.
\end{proof}

\section{Morphisms}
To establish a category of varieties we define morphisms as continuous maps preserving regular functions. The central principle is that affine varieties are controlled by their coordinate rings.

\subsection{Regular functions}
\begin{defn}\label{def:regular_function_quasiaffine}
Let $X$ be a quasi-affine variety. A function $f:X\to k$ is \textbf{regular at} $P\in X$ if there is an open neighbourhood $U$ of $P$ and polynomials $g,h$ with $h$ nowhere zero on $U$ such that $f|_U=g/h$. The function is \textbf{regular} if it is regular at every point of $X$.
\end{defn}
On an irreducible variety, every regular function on a non-empty open set determines a rational function. What may fail is the existence of a single quotient expression valid on every affine chart at once.

\begin{defn}
Let $Y$ be a variety. The \textbf{function field} $K(Y)$ is the field of equivalence classes of pairs $(U,f)$, where $U\subseteq Y$ is non-empty open and $f$ is regular on $U$, modulo agreement on some non-empty open subset of the intersection of domains.
\end{defn}

\begin{proposition}\label{prop:affine_basics}
Let $X$ be an affine variety. Then
\begin{enumerate}
\item $A(X)\cong\mathcal O(X)$;
\item for $P=(P_1,\ldots,P_n)\in X$,
\[
\mathcal O_{P,X}\cong A(X)_{\mathfrak m_P},\qquad \mathfrak m_P=(x_1-P_1,\ldots,x_n-P_n);
\]
\item $K(X)\cong\operatorname{Frac} A(X)$, and $K(X)/k$ is a finitely generated field extension with $\operatorname{tr.deg}_kK(X)=\dim X$.
\end{enumerate}
\end{proposition}
\begin{proof}
The natural map $A(X)\to\mathcal O(X)$ is injective by the definition of $I(X)$. A regular function is locally a quotient of polynomial functions, so it lies in every localisation $A(X)_{\mathfrak m}$ inside $\operatorname{Frac} A(X)$. Since $A(X)$ is a finitely generated domain, $A(X)=\bigcap_{\mathfrak m\in\operatorname{Max}A(X)}A(X)_{\mathfrak m}$ inside $\operatorname{Frac} A(X)$, proving global regular functions are polynomial. The local-ring and function-field statements follow from the same local quotient description. The transcendence-degree equality is Theorem \ref{thm:onepointeight}.
\end{proof}

\subsection{Morphisms}
\begin{defn}
A \textbf{variety} is an affine, quasi-affine, projective, or quasi-projective variety. A \textbf{morphism} $f:X\to Y$ is a continuous map such that regular functions on open subsets of $Y$ pull back to regular functions on the corresponding inverse-image open subsets of $X$.
\end{defn}
\begin{lemma}\label{lem:check_variables}
Let $X$ be a variety and let $Y\subseteq\bb A^n$ be affine. A map $f:X\to Y$ is a morphism if and only if each coordinate function $x_i\circ f$ is regular on $X$.
\end{lemma}
\begin{proof}
One direction is immediate. Conversely, polynomial functions in the coordinates pull back to regular functions, and regular functions on affine open subsets are locally quotients of such polynomial functions.
\end{proof}
\begin{proposition}\label{prop:adjunction}
Let $X$ be any variety and let $Y$ be affine. There is a natural bijection
\[
\Hom_{\mathbf{Var}}(X,Y)\cong\Hom_{k\operatorname{-alg}}(A(Y),\mathcal O(X)).
\]
Under this bijection, $f:X\to Y$ is dominant if and only if $f^*:A(Y)\to\mathcal O(X)$ is injective.
\end{proposition}
\begin{proof}
A morphism gives pullback of regular functions. Conversely, a homomorphism $\phi:A(Y)\to\mathcal O(X)$ sends the coordinate classes $\bar x_i$ to regular functions on $X$, defining $P\mapsto(\phi(\bar x_i)(P))$. The polynomial relations in $I(Y)$ map to zero, so the image lies in $Y$, and Lemma \ref{lem:check_variables} proves this is a morphism. The dominance statement says exactly that no nonzero regular function on $Y$ vanishes on the image.
\end{proof}

\subsection{Products}
The product of varieties is not generally the product topological space. Unless explicitly stated otherwise, $X\times Y$ denotes the categorical product variety.
\begin{proposition}\label{prop:prod_affine_variety}
If $X\subseteq\bb A^n$ and $Y\subseteq\bb A^m$ are affine varieties, then $X\times Y\subseteq\bb A^{n+m}$ is affine and
\[
A(X\times Y)\cong A(X)\otimes_kA(Y).
\]
It is the categorical product in varieties.
\end{proposition}
\begin{proof}
Writing $A(X)=k[x]/I$ and $A(Y)=k[y]/J$, one has $A(X\times Y)=k[x,y]/(I+J)\cong A(X)\otimes_kA(Y)$. The universal property follows from Proposition \ref{prop:adjunction}.
\end{proof}
\begin{proposition}\label{prop:prod_projective_variety}
If $X\subseteq\bb P^r$ and $Y\subseteq\bb P^s$ are projective varieties, their product is realised inside $\bb P^{(r+1)(s+1)-1}$ by the Segre embedding.
\end{proposition}
\begin{proof}
Use Proposition \ref{prop:segre_isomorphism}; the universal property is checked on standard affine charts, reducing to Proposition \ref{prop:prod_affine_variety}.
\end{proof}

\subsection{Rational maps}
\begin{defn}
A \textbf{rational map} $X\dashrightarrow Y$ is an equivalence class of pairs $(U,\varphi)$ where $U\subseteq X$ is non-empty open and $\varphi:U\to Y$ is a morphism. Two representatives are equivalent if they agree on some non-empty open subset of the intersection of their domains. A rational map is \textbf{dominant} if one, hence every, representative has dense image.
\end{defn}
\begin{lemma}\label{lem:ambient_hypersurface_affine}
If $H=Z(f)\subseteq\bb A^n$ is a hypersurface, then $\bb A^n\setminus H$ is affine.
\end{lemma}
\begin{proof}
It is isomorphic to $Z(tf-1)\subseteq\bb A^{n+1}$ by projection to the first $n$ coordinates, with inverse $P\mapsto(P,1/f(P))$.
\end{proof}
\begin{proposition}\label{prop:base}
Open affine subvarieties form a basis for the topology of every variety.
\end{proposition}
\begin{proof}
For a quasi-affine $Y\subseteq\bb A^n$, a point $P\in Y$, and an open neighbourhood $U$, choose $f$ vanishing on $\overline Y\setminus U$ but not at $P$. Then $Y\cap D(f)$ is contained in $U$ and is closed in the affine variety $D(f)\cong Z(tf-1)$, hence affine. The projective and quasi-projective cases reduce to affine charts.
\end{proof}
\begin{thm}\label{thm:equivalence_of_cats}
The category of varieties with dominant rational maps is anti-equivalent to the category of finitely generated field extensions of $k$; equivalently,
\[
\Hom_{\mathrm{dom.rat}}(X,Y)\cong\Hom_{k\operatorname{-alg}}(K(Y),K(X)).
\]
\end{thm}
\begin{proof}
A dominant rational map pulls rational functions on $Y$ back to rational functions on $X$. Conversely, a $k$-embedding $K(Y)\hookrightarrow K(X)$ sends the coordinate functions of an affine open $Y'\subseteq Y$ to rational functions on $X$; after shrinking to a common domain of regularity, these define a morphism from a non-empty open subset of $X$ to $Y'$. The two constructions are inverse up to the equivalence relation on rational maps.
\end{proof}
\begin{cor}\label{cor:function_field_birational}
For varieties $X,Y$, the following are equivalent: $X$ and $Y$ are birational; non-empty open subsets of $X$ and $Y$ are isomorphic; and $K(X)\cong K(Y)$ over $k$.
\end{cor}
\begin{proof}
This is the isomorphism case of Theorem \ref{thm:equivalence_of_cats}.
\end{proof}
\begin{proposition}\label{prop:hypersurface}
Every variety $X$ of dimension $r$ is birationally equivalent to a hypersurface in $\bb A^{r+1}$ and, after projective closure, to a hypersurface in $\bb P^{r+1}$.
\end{proposition}
\begin{proof}
Choose a separating transcendence basis $x_1,\ldots,x_r\in K(X)$. Since $k$ is algebraically closed and hence perfect, $K(X)$ is finite separable over $k(x_1,\ldots,x_r)$, so by the primitive element theorem $K(X)=k(x_1,\ldots,x_r,y)$ for some $y\in K(X)$. Clearing denominators in the irreducible equation of $y$ gives a hypersurface with function field $K(X)$.
\end{proof}

\section{Singularities}
\subsection{Singular points}
Throughout this subsection, $X\subseteq\bb A^n$ is an affine variety of dimension $r$.
\begin{defn}\label{def:singular_point}
Let $I(X)=(f_1,\ldots,f_m)$. For $P\in X$, the \textbf{Jacobian matrix} at $P$ is
\[
J_X(P)=\left(\frac{\partial f_i}{\partial x_j}(P)\right)_{1\leq i\leq m,\,1\leq j\leq n}.
\]
The point $P$ is \textbf{nonsingular} if $\operatorname{rank}J_X(P)=n-r$, and \textbf{singular} otherwise.
\end{defn}
\begin{remark}\label{rmk:defining_functions_invariant}
The rank condition in Definition \ref{def:singular_point} is independent of the chosen generators of $I(X)$. If $f_i=\sum_j h_{ij}g_j$, then at a point of $X$ the rows of the Jacobian matrix of the $f_i$ are linear combinations of the rows of the Jacobian matrix of the $g_j$. Reversing the roles of the two generating sets gives equality of ranks.
\end{remark}
\begin{defn}\label{def:regular_local}
Let $(A,\mathfrak m)$ be a Noetherian local ring with residue field $\kappa=A/\mathfrak m$. The ring $A$ is \textbf{regular} if
\[
\dim A=\dim_\kappa \mathfrak m/\mathfrak m^2.
\]
A variety $Y$ is \textbf{nonsingular at} $P$ if the local ring $\mathcal O_{Y,P}$ is regular.
\end{defn}
\begin{lemma}\label{lem:equivalence_singularity}
For an affine variety $Y\subseteq\bb A^n$, the Jacobian definition of nonsingularity agrees with regularity of the local ring $\mathcal O_{Y,P}$.
\end{lemma}
\begin{proof}
Let $A=k[x_1,\ldots,x_n]$, let $\mathfrak a=I(Y)$, and let $\mathfrak m_P=(x_1-P_1,\ldots,x_n-P_n)$. The differential map identifies $\mathfrak m_P/\mathfrak m_P^2$ with $k^n$, and the subspace $(\mathfrak a+\mathfrak m_P^2)/\mathfrak m_P^2$ is spanned by the gradients of generators of $\mathfrak a$ at $P$. Hence
\[
\dim_k \mathfrak m_{Y,P}/\mathfrak m_{Y,P}^2=n-\operatorname{rank}J_Y(P).
\]
Since $\dim\mathcal O_{Y,P}=r$ at a closed point of an irreducible affine variety, regularity is equivalent to $n-\operatorname{rank}J_Y(P)=r$, i.e. to $\operatorname{rank}J_Y(P)=n-r$.
\end{proof}
\begin{lemma}\label{lem:projective_jacobian}
Let $X=Z(f_1,\ldots,f_m)\subseteq\bb P^n$ be a projective variety of dimension $r$, and let $P\in X$. Then $X$ is nonsingular at $P$ if and only if the homogeneous Jacobian matrix
\[
\left(\frac{\partial f_i}{\partial x_j}(P)\right)_{i,j}
\]
has rank $n-r$.
\end{lemma}
\begin{proof}
Choose a chart $U_l$ containing $P$ and dehomogenise by setting $x_l=1$. The affine Jacobian criterion says that nonsingularity is equivalent to rank $n-r$ for the dehomogenised equations. Euler's formula for homogeneous polynomials gives
\[
\sum_{j=0}^n P_j\frac{\partial f_i}{\partial x_j}(P)=\deg(f_i)f_i(P)=0,
\]
so the radial direction gives the one homogeneous column relation removed on passing to the affine chart. Thus the affine and homogeneous rank conditions agree.
\end{proof}
\begin{thm}\label{thm:singular_locus_closed}
For any variety $Y$, the singular locus $\operatorname{Sing}Y$ is a proper closed subset of $Y$.
\end{thm}
\begin{proof}
Closedness is local on an affine open cover. If $Y=Z(f_1,\ldots,f_m)\subseteq\bb A^n$ has dimension $r$, then $\operatorname{Sing}Y$ is cut out by $I(Y)$ together with all $(n-r)\times(n-r)$ minors of the Jacobian matrix. Thus it is closed. Properness follows because over the algebraically closed field $k$, hence over a perfect field, the nonsingular locus of an irreducible variety is nonempty; equivalently, apply the Jacobian criterion together with generic smoothness on an affine open subset.
\end{proof}

\section{Blowups}
The blowup of the origin in affine space records limiting directions of lines approaching the origin.
\begin{defn}\label{def:general_blowup}
The blowup of $\bb A^n$ at the origin is
\[
\operatorname{Bl}_0\bb A^n=\bigl\{(x,[P])\in\bb A^n\times\bb P^{n-1}:x_iP_j=x_jP_i\text{ for all }i,j\bigr\}.
\]
The projection $\pi:\operatorname{Bl}_0\bb A^n\to\bb A^n$ is induced by projection onto the first factor. If $Y\subseteq\bb A^n$ is an affine variety passing through the origin, the \textbf{blowup of $Y$ at the origin} is the strict transform
\[
\operatorname{Bl}_0Y=\overline{\pi^{-1}(Y\setminus\{0\})}\subseteq\operatorname{Bl}_0\bb A^n.
\]
\end{defn}
For plane curves, the two affine charts of $\bb P^1$ give the equations $y=zx$ and $x=zy$. Both charts are needed, since a single chart can miss vertical tangent directions.
\begin{defn}\label{def:plane_curve_blowup}
Let $Y=Z(f)\subseteq\bb A^2$ be a plane curve through the origin. In the chart $y=zx$, substitute $y=zx$ and write
\[
f(x,zx)=x^m\bar f(x,z),\qquad \bar f(0,z)\not\equiv0,
\]
where $m$ is the multiplicity of $f$ at the origin. The strict transform in this chart is
\[
Z(\bar f(x,z),\,y-zx)\subseteq\bb A^3.
\]
The second chart is obtained analogously by substituting $x=zy$.
\end{defn}
\begin{example}\label{ex:node_blowup}
Let $Y=Z(y^2-x^2(x+1))$. The origin is singular because the gradient vanishes there. In the chart $y=zx$ one gets
\[
z^2x^2-x^2(x+1)=x^2(z^2-x-1),
\]
so the strict transform is cut out by $z^2-x-1$ together with $y=zx$. Its Jacobian has full rank everywhere on this chart, so this chart of the strict transform is nonsingular.
\end{example}
\begin{defn}
Let $Y=Z(f)$ be a plane curve and $P\in Y$. After translating $P$ to the origin, write
\[
f=f_m+f_{m+1}+\cdots
\]
where $f_i$ is homogeneous of degree $i$ and $f_m\neq0$. The integer $m$ is the \textbf{multiplicity} of $Y$ at $P$, and the linear factors of the lowest nonzero homogeneous part $f_m$ are the \textbf{tangent directions}. A singular point of multiplicity $2$ with distinct tangent directions is a \textbf{node}.
\end{defn}
\clearpage
\section{Hartshorne Exercise Solutions, Corrected}
\label{sec:hartshorne-solutions}
The exercise statements below are working paraphrases of the corresponding problems in Hartshorne's \emph{Algebraic Geometry} \cite{hartshorne}. They are written so that each solution can be read independently, without reproducing the official text verbatim.

\newcommand{\exhead}[1]{\subsubsection*{Exercise #1}\addcontentsline{toc}{subsubsection}{Exercise #1}}
\newcommand{\questionpar}{\paragraph{Question, paraphrased.}}
\newcommand{\solutionpar}{\paragraph{Solution.}}

\subsection{Chapter I, \S 1: Affine varieties}

\exhead{I.1.1}
\questionpar Classify the standard affine curves appearing in the first exercise by computing their coordinate rings: the parabola $y=x^2$, the hyperbola $xy=1$, and an irreducible affine conic.
\solutionpar For $Y=Z(y-x^2)\subseteq \bb A^2$,
\[
A(Y)=k[x,y]/(y-x^2)\cong k[x].
\]
For $Z=Z(1-xy)$,
\[
A(Z)=k[x,y]/(1-xy)\cong k[x,x^{-1}],
\]
via $y\mapsto x^{-1}$. If $f(x,y)$ is an irreducible quadratic and $\operatorname{char}k\neq2$, an affine linear change of variables brings its quadratic part either to a square or to a product of independent linear forms. In the first case the curve is isomorphic to $y=x^2$; in the second it is isomorphic to $xy=1$ after completing the square or translating variables. In characteristic $2$ one should replace this diagonalisation argument by the usual classification of nonsingular conics.

\exhead{I.1.2}
\questionpar Study the affine curve $Y\subseteq\bb A^3$ parametrised by $(t,t^2,t^3)$. Find equations, show it is isomorphic to $\bb A^1$, and compute its ideal.
\solutionpar The image is contained in
\[
Z(y-x^2,z-x^3).
\]
The map $k[x,y,z]\to k[t]$ sending $x\mapsto t$, $y\mapsto t^2$, $z\mapsto t^3$ is surjective and has kernel $(y-x^2,z-x^3)$. Hence
\[
A(Y)\cong k[t],
\]
so $Y\cong\bb A^1$ and $I(Y)=(y-x^2,z-x^3)$.

\exhead{I.1.3}
\questionpar Decompose the algebraic set in $\bb A^3$ defined by $x^2-yz$ and $xz-x$ into irreducible components and identify their prime ideals.
\solutionpar The equations are
\[
x^2=yz,
\qquad x(z-1)=0.
\]
If $x=0$, then $yz=0$, giving the two lines $Z(x,y)$ and $Z(x,z)$. If $z=1$, then $y=x^2$, giving the parabola $Z(z-1,y-x^2)$. Thus
\[
Y=Z(x,y)\cup Z(x,z)\cup Z(z-1,y-x^2),
\]
with prime ideals $(x,y)$, $(x,z)$, and $(z-1,y-x^2)$.

\exhead{I.1.4}
\questionpar Compare the Zariski topology on $\bb A^2$ with the product of the Zariski topologies on two copies of $\bb A^1$.
\solutionpar The diagonal
\[
\Delta=\{(t,t):t\in k\}\subseteq\bb A^1\times\bb A^1
\]
is Zariski closed in $\bb A^2$, since $\Delta=Z(y-x)$. In the product topology, closed sets are finite unions of products of closed subsets of $\bb A^1$. Since proper closed subsets of $\bb A^1$ are finite, an infinite proper closed set such as $\Delta$ cannot be closed in the product topology. Thus the Zariski topology on $\bb A^2$ is strictly finer.

\exhead{I.1.5}
\questionpar Characterise the finitely generated reduced $k$-algebras as coordinate rings of affine algebraic sets.
\solutionpar If $B$ is finitely generated and reduced, write $B\cong k[x_1,\ldots,x_n]/\mathfrak a$. Since $B$ is reduced, $\mathfrak a$ is radical. By the Nullstellensatz,
\[
A(Z(\mathfrak a))=k[x_1,\ldots,x_n]/I(Z(\mathfrak a))=k[x_1,\ldots,x_n]/\sqrt{\mathfrak a}\cong B.
\]
Conversely, coordinate rings of affine algebraic sets are finitely generated and reduced.

\exhead{I.1.6}
\questionpar Show that a nonempty open subset of an irreducible space is irreducible and dense.
\solutionpar Let $X$ be irreducible and $U\subseteq X$ nonempty open. Any two nonempty open subsets of $U$ are intersections of $U$ with nonempty open subsets of $X$, hence meet. Thus $U$ is irreducible. Also every nonempty open subset of an irreducible space is dense, because its complement cannot contain a nonempty open subset.

\exhead{I.1.8}
\questionpar If $Y\subseteq\bb A^n$ is irreducible of dimension $r$ and $H$ is a hypersurface not containing $Y$, compute the dimensions of the irreducible components of $Y\cap H$.
\solutionpar Let $H=Z(f)$ and pass to the domain $A(Y)$. Since $H$ does not contain $Y$, the image of $f$ is a nonzero nonunit of $A(Y)$. The irreducible components of $Y\cap H$ correspond to minimal primes over $(f)$. By Krull's principal ideal theorem, these minimal primes have height at most $1$, and since $A(Y)$ is a domain and $f\neq0$, their height is exactly $1$. Therefore each component has dimension $r-1$.

\exhead{I.1.9}
\questionpar If an ideal in $k[x_1,\ldots,x_n]$ can be generated by $r$ elements, show that every irreducible component of its zero set has dimension at least $n-r$.
\solutionpar Let $\mathfrak p$ be a minimal prime over $(f_1,\ldots,f_r)$. By the generalized principal ideal theorem, $\height\mathfrak p\leq r$. Therefore
\[
\dim k[x_1,\ldots,x_n]/\mathfrak p=n-\height\mathfrak p\geq n-r.
\]
This is the dimension of the corresponding irreducible component.

\exhead{I.1.10}
\questionpar Prove basic properties of topological dimension: monotonicity for subspaces, equality for finite open covers, examples with dense opens, and behaviour under irreducibility assumptions.
\solutionpar Monotonicity follows by viewing a chain of irreducible closed subsets of a subspace as a chain of closures in the ambient space, when the subspace is closed. For a finite open cover of an irreducible Noetherian space, Lemma \ref{lem:the_lemma_never_stated} identifies chains in an open set with chains meeting that open set. Thus dimension is the supremum of the dimensions of the members of the cover. A dense open subset of an irreducible variety has the same dimension as the variety, but arbitrary dense open subsets of non-Noetherian spaces need not behave this way.

\exhead{I.1.12}
\questionpar Give an example over $\bb R$ of an irreducible polynomial whose zero set is not irreducible.
\solutionpar The polynomial $x^2+y^2+1\in\bb R[x,y]$ is irreducible, but its real zero set is empty. By the convention that the empty set is not irreducible, this gives the desired example.

\subsection{Chapter I, \S 2: Projective varieties}

\exhead{I.2.2}
\questionpar For a homogeneous ideal $\mathfrak a\subseteq S=k[x_0,\ldots,x_n]$, prove the standard equivalent criteria for $Z(\mathfrak a)=\varnothing$ in projective space.
\solutionpar The correct statement is that $Z(\mathfrak a)=\varnothing$ if and only if the radical of $\mathfrak a$ contains the irrelevant ideal $S_+=(x_0,\ldots,x_n)$, equivalently if and only if $S_d\subseteq\mathfrak a$ for all sufficiently large $d$. One must not say that the pure powers $x_i^d$ generate $S_d$ as a vector space. Rather, if $x_i^N\in\mathfrak a$ for all $i$, then every monomial of degree $(n+1)N$ is divisible by some $x_i^N$, so $S_d\subseteq\mathfrak a$ for all sufficiently large $d$.

\exhead{I.2.3}
\questionpar Establish the projective versions of the order-reversing identities between homogeneous ideals and projective algebraic sets.
\solutionpar The usual inclusions $T_1\subseteq T_2\Rightarrow Z(T_1)\supseteq Z(T_2)$ and $Y_1\subseteq Y_2\Rightarrow I(Y_1)\supseteq I(Y_2)$ hold as in the affine case. The projective Nullstellensatz gives, for homogeneous ideals,
\[
I(Z(\mathfrak a))=\{f\in S\text{ homogeneous}: f^N\in\mathfrak a\text{ for some }N>0\},
\]
with the standard empty-set exception involving the irrelevant ideal. The radical of a homogeneous ideal is homogeneous, meaning generated by homogeneous elements, not that every element of the radical is homogeneous.

\exhead{I.2.6}
\questionpar Relate the dimension of a projective variety $Y$ to the dimension of its homogeneous coordinate ring.
\solutionpar Cover $Y\subseteq\bb P^n$ by the standard affine opens $U_i$. On $U_i$, the coordinate ring is the degree-zero part of the localisation $S(Y)_{x_i}$. Its dimension is one less than $\dim S(Y)$, because passing to the affine cone adds the radial coordinate. Thus
\[
\dim S(Y)=\dim Y+1.
\]

\exhead{I.2.9}
\questionpar For the twisted cubic in $\bb P^3$, compute the homogeneous ideal of its projective closure.
\solutionpar Let $[w:x:y:z]$ be coordinates and parametrise the twisted cubic by
\[
[s:t]\longmapsto [s^3:s^2t:st^2:t^3].
\]
The ideal is
\[
I(C)=(wy-x^2,\; wz-xy,\; xz-y^2).
\]
These three quadrics vanish on the parametrisation. Conversely, on the chart $w\neq0$ they give $y=x^2/w$ and $z=x^3/w^2$, recovering the affine twisted cubic; the point at infinity is obtained when $w=0$. A standard Gr\"obner basis or Hilbert polynomial check shows that these equations define no extra projective component.

\exhead{I.2.11}
\questionpar Characterise linear subvarieties of projective space.
\solutionpar A projective linear subvariety of $\bb P^n$ is the zero set of a vector subspace of the dual space of linear forms. If it has codimension $m$, then it can be cut out by $m$ linearly independent linear forms. The independence is essential: without it, the dimension count fails.

\exhead{I.2.12}
\questionpar Prove the basic properties of the $d$-uple Veronese embedding.
\solutionpar The Veronese map sends
\[
[x_0:\cdots:x_n]\longmapsto [\text{all monomials of degree }d].
\]
Its image is closed because its homogeneous coordinate ring is the subring generated by degree-$d$ monomials, and the kernel of the corresponding polynomial map gives homogeneous equations for the image. It is injective: if two points have the same degree-$d$ monomial coordinates, choose $i$ with $P_i\neq0$. The coordinate $x_i^d$ is nonzero at both points, so on the affine chart $U_i$ the ratios $x_j/x_i$ are recovered from $x_jx_i^{d-1}/x_i^d$. The same chartwise formula gives a regular inverse onto the image, so the map is an isomorphism onto a projective variety.

\exhead{I.2.14}
\questionpar Prove that the Segre map is an isomorphism onto its image and identify its equations.
\solutionpar The Segre image in $\bb P^{(r+1)(s+1)-1}$ is cut out by the $2\times2$ minors
\[
z_{ij}z_{kl}-z_{il}z_{kj}.
\]
If $P=[P_{ij}]$ satisfies these minors and $P_{ab}\neq0$, then
\[
P_{ij}P_{ab}=P_{ib}P_{aj},
\]
so $P$ equals the Segre image of $[P_{0b}:\cdots:P_{rb}]$ and $[P_{a0}:\cdots:P_{as}]$. This avoids division by possibly zero entries such as $P_{kb}$.

\exhead{I.2.16}
\questionpar Give examples showing that intersections of varieties need not be varieties and that ideals of intersections need not be sums of ideals.
\solutionpar Two irreducible varieties may have reducible intersection: for instance, suitable quadric surfaces in $\bb P^3$ may intersect in a union of curves. For the twisted cubic $C\subseteq\bb P^3$, the two quadrics
\[
wy-x^2=0,
\qquad wz-xy=0
\]
cut out $C$ together with the residual line $w=x=0$. Scheme-theoretically this is reflected by the saturated decomposition
\[
(wy-x^2,wz-xy)=(I(C))\cap(w,x)
\]
up to the usual saturation check.

\subsection{Chapter I, \S 3: Morphisms}

\exhead{I.3.1}
\questionpar Classify nonsingular affine and projective conics up to isomorphism.
\solutionpar Over an algebraically closed field of characteristic not $2$, a nonsingular projective plane conic can be diagonalised and is isomorphic to $\bb P^1$. Affine conics are obtained by deleting the hyperplane section at infinity. Depending on that intersection, one obtains the affine line, the punctured affine line, or the standard affine conic models appearing in Exercise I.1.1. In characteristic $2$, one should use the general theorem that every nonsingular plane conic over an algebraically closed field is isomorphic to $\bb P^1$, rather than an orthogonal diagonalisation argument.

\exhead{I.3.2}
\questionpar Analyse the cusp $y^2=x^3$ and related rational curves by parametrisation.
\solutionpar The cusp is parametrised by
\[
x=t^2,
\qquad y=t^3,
\]
since $y^2=t^6=x^3$. The induced map $k[x,y]/(y^2-x^3)\to k[t]$ is injective but not surjective, because $t$ is not in the image. Thus the curve is rational but not isomorphic to $\bb A^1$.

\exhead{I.3.5}
\questionpar Show that the complement of a hyperplane in projective space is affine, and use projective linear changes of coordinates correctly.
\solutionpar A projective linear automorphism carries any hyperplane to a coordinate hyperplane. Thus it suffices to consider $H=Z(x_0)\subseteq\bb P^n$. The complement $\bb P^n\setminus H$ is the standard affine chart $U_0\cong\bb A^n$.

\exhead{I.3.6}
\questionpar Show that $\bb A^2\setminus\{0\}$ is not affine.
\solutionpar Let $X=\bb A^2\setminus\{0\}$. A regular function on $X$ is locally a quotient of polynomials. Since the missing set has codimension $2$, the same argument as Hartogs' phenomenon in this algebraic setting gives
\[
\Gamma(X,\mathcal O_X)=k[x,y].
\]
The inclusion $j:X\hookrightarrow\bb A^2$ induces an isomorphism on global regular functions. If $X$ were affine, the affine coordinate-ring equivalence would force $j$ to be an isomorphism. But $j$ is not surjective, contradiction.

\exhead{I.3.8}
\questionpar Show that regular functions on a projective variety are constant.
\solutionpar It suffices to prove this for an irreducible projective variety $Y$. A regular function $f:Y\to\bb A^1$ gives a morphism from a projective variety to an affine variety. The image is both projective and affine; equivalently, one can use homogeneous coordinates to see that a globally regular quotient of homogeneous forms of the same degree has no poles only when it is constant. Hence $\mathcal O(Y)=k$.

\exhead{I.3.15}
\questionpar Prove the product properties for affine varieties, including the formula for coordinate rings and dimensions.
\solutionpar If $X$ and $Y$ are affine, then
\[
A(X\times Y)\cong A(X)\otimes_k A(Y).
\]
Since $k$ is algebraically closed and $A(X),A(Y)$ are finitely generated $k$-domains, the tensor product is again a domain. Therefore
\[
\dim(X\times Y)=\operatorname{trdeg}_k(A(X)\otimes_k A(Y))=\dim X+
\dim Y.
\]
The universal property of products follows from the affine morphism-coordinate-ring correspondence.

\exhead{I.3.17}
\questionpar Check normality for standard examples: affine space, a quadric surface, and a cusp.
\solutionpar Affine space is normal because $k[x_1,\ldots,x_n]$ is a UFD. The quadric surface $Q=Z(xw-yz)\subseteq\bb A^4$ is normal by Serre's criterion: it is a hypersurface, hence Cohen-Macaulay, and its singular locus has codimension at least $2$ in the cases under consideration. The cusp $A=k[t^2,t^3]$ is not normal, since $t$ is integral over $A$ by the equation $T^2-t^2=0$ but $t\notin A$.

\subsection{Chapter I, \S 4: Rational maps}

\exhead{I.4.1}
\questionpar Show that regular functions glue over two open subsets.
\solutionpar Let $U,V\subseteq X$ be open and let $f\in\mathcal O(U)$, $g\in\mathcal O(V)$ agree on $U\cap V$. Define $h$ by $h=f$ on $U$ and $h=g$ on $V$. Regularity is local, so every point of $U\cup V$ has a neighbourhood on which $h$ is given by a quotient of polynomials. Thus $h\in\mathcal O(U\cup V)$.

\exhead{I.4.2}
\questionpar Prove the analogous gluing statement for rational maps.
\solutionpar If morphisms $\varphi_i:U_i\to Y$ agree on overlaps, then for every regular function $a$ on an open subset of $Y$, the pullbacks $a\circ\varphi_i$ glue by Exercise I.4.1. Hence the glued set-map is a morphism. Rational maps are then glued by restricting to common nonempty open domains.

\exhead{I.4.4}
\questionpar Characterise rational curves and give standard examples.
\solutionpar A curve is rational if its function field is $k(t)$, equivalently if it is birational to $\bb P^1$. The affine line, the projective line, and any nonsingular conic are rational. The cusp $y^2=x^3$ is rational via $x=t^2$, $y=t^3$, but this parametrisation is not an isomorphism at the cusp.

\subsection{Chapter I, \S 5: Nonsingular varieties}

\exhead{I.5.9}
\questionpar Show that a nonsingular plane projective curve is irreducible, with attention to repeated factors.
\solutionpar Let $C=Z(f)\subseteq\bb P^2$. If $f=gh$ with $g,h$ nonconstant and coprime, then $Z(g)$ and $Z(h)$ meet in $\bb P^2$, so at an intersection point all partial derivatives of $f=gh$ vanish, giving a singular point. If $f$ has a repeated factor, say $f=g^2h$, then along $Z(g)$ the partial derivatives of $f$ vanish as well. Therefore a nonsingular plane projective curve is defined by an irreducible reduced homogeneous polynomial.

\subsection{Chapter II, \S 1: Sheaves}

\exhead{II.1.1}
\questionpar Describe the constant sheaf associated to an abelian group $A$.
\solutionpar The constant presheaf $U\mapsto A$ need not be a sheaf. Its sheafification is the sheaf of locally constant $A$-valued functions. Thus
\[
\mathscr A(U)=\{\text{locally constant functions }U\to A\}.
\]
If the connected components of $U$ are open, this is $\prod_{\pi_0(U)}A$. In particular, if $U$ is connected, then $\mathscr A(U)\cong A$.

\subsection{Chapter II, \S 2: Schemes}

\exhead{II.2.3}
\questionpar Construct the reduced scheme associated to a scheme $X$.
\solutionpar On an affine open $U=\Spec R$, define
\[
U_{\mathrm{red}}=\Spec(R/\sqrt0).
\]
These affine constructions glue because nilradicals commute with localisation. The underlying topological space is the same as that of $X$, and the structure sheaf is
\[
\mathcal O_{X_{\mathrm{red}}}=\mathcal O_X/\mathcal N,
\]
where $\mathcal N$ is the sheaf of nilpotent elements. The natural closed immersion $X_{\mathrm{red}}\to X$ is universal for morphisms from reduced schemes to $X$.

\exhead{II.2.8}
\questionpar Interpret $k[\varepsilon]/(\varepsilon^2)$-valued points as tangent vectors.
\solutionpar A $k[\varepsilon]/(\varepsilon^2)$-valued point of a $k$-scheme $X$ supported at a $k$-rational point $x$ is equivalent to a local $k$-algebra map
\[
\mathcal O_{X,x}\to k[\varepsilon]/(\varepsilon^2)
\]
whose composition with the quotient to $k$ is evaluation at $x$. Such a map has the form
\[
a\longmapsto \bar a+D(a)\varepsilon,
\]
where $D$ is a $k$-derivation $\mathcal O_{X,x}\to k$. Since $D$ vanishes on $\mathfrak m_x^2$, this is the same as an element of
\[
\operatorname{Hom}_k(\mathfrak m_x/\mathfrak m_x^2,k)=T_xX.
\]

\exhead{II.2.8, generic points}
\questionpar Show that irreducible closed subsets of a scheme have unique generic points.
\solutionpar Affine schemes are sober: an irreducible closed subset of $\Spec A$ is $V(\mathfrak p)$ for a unique prime $\mathfrak p$, and $\mathfrak p$ is its generic point. For a general scheme $X$, let $Z\subseteq X$ be irreducible closed and choose an affine open $U$ meeting $Z$. Then $Z\cap U$ is irreducible closed in $U$, so it has a generic point $\xi$. Since $Z\cap U$ is dense in $Z$, the closure of $\xi$ in $X$ is $Z$. Uniqueness follows by restricting to affine open subsets.

\exhead{II.2.11}
\questionpar Describe $\Spec\mathbb F_p[x]$ and the residue fields at its points.
\solutionpar The prime ideals of $\mathbb F_p[x]$ are
\[
(0)\quad\text{and}\quad(f),
\]
where $f$ is irreducible. The point $(0)$ is generic. Each nonzero prime $(f)$ is closed, with residue field
\[
\kappa((f))=\mathbb F_p[x]/(f)\cong\mathbb F_{p^{\deg f}}.
\]
Thus there are closed points of every finite residue degree.

\exhead{II.2.16(c)}
\questionpar Extend a section over a principal open after multiplying by a power of the defining function.
\solutionpar If $X$ is affine, $X=\Spec A$, and $X_f=D(f)$, then
\[
\Gamma(D(f),\mathcal O_X)=A_f.
\]
Thus any section $b$ over $D(f)$ can be written as $a/f^N$ for some $a\in A$ and $N\geq0$. Equivalently, $f^Nb$ is the restriction of a global section. The conclusion is not generally that $b$ itself is the restriction of a global section.

\subsection{Chapter II, \S 3: First properties of schemes}

\exhead{II.3.1}
\questionpar Show that being locally of finite type is local on the target.
\solutionpar Let $f:X\to Y$ be a morphism. If $Y$ has an affine open cover $U_i=\Spec A_i$ such that $f^{-1}(U_i)$ is covered by affine opens $\Spec B_{ij}$ with each $B_{ij}$ a finitely generated $A_i$-algebra, then the same property holds after refining to principal affine opens. The algebra input is: if $B$ is an $A_f$-algebra, then $B$ is finitely generated over $A_f$ if and only if it is finitely generated over $A$ after remembering that $A_f=A[1/f]$ is itself finitely generated as an $A$-algebra. Refining and gluing these affine statements proves locality on the target.

\exhead{II.3.14}
\questionpar Prove that closed points are dense in a scheme of finite type over a field.
\solutionpar Cover $X$ by affine opens $\Spec A_i$ with each $A_i$ a finitely generated $k$-algebra. By the Jacobson property for finitely generated algebras over a field, closed points are dense in each nonempty affine open. If a point $x$ has residue field finite over $k$, then its closure has function field $\kappa(x)$, hence dimension
\[
\operatorname{trdeg}_k\kappa(x)=0,
\]
so $x$ is closed in $X$. Therefore every nonempty open subset contains a closed point.

\subsection{B\'ezout's theorem for plane curves}

\exhead{B\'ezout}
\questionpar Prove the basic form of B\'ezout's theorem for two plane projective curves with no common component.
\solutionpar Let $F,G\in S=k[x,y,z]$ be homogeneous of positive degrees $m,n$ with no common factor. Since $S$ is a UFD, $F,G$ form a homogeneous regular sequence. Therefore
\[
H_{S/(F,G)}(t)=\frac{(1-t^m)(1-t^n)}{(1-t)^3}.
\]
For large $d$, the Hilbert function of $S/(F,G)$ is constant and equal to $mn$. This constant is the length of the zero-dimensional scheme $Z(F,G)\subseteq\bb P^2$, counted with scheme-theoretic multiplicity. Hence the sum of intersection multiplicities is $mn$.


\clearpage
\section{Errata incorporated in this combined version}
\label{sec:errata-incorporated}
This section records the main changes made while combining the original varieties notes with the Hartshorne solutions.
\begin{enumerate}
\item The convention on varieties was made explicit: affine and projective varieties are irreducible, while algebraic sets need not be.
\item The Nullstellensatz notation was corrected from $(X)$ to $(\Phi)$, and the closure identity $ZI(T)=\overline T$ was reproved cleanly.
\item The false irreducibility-transfer lemma was replaced by the correct statement that irreducibility of a subspace is intrinsic to its induced topology.
\item The discussion of linear affine varieties was corrected: zero sets of linear polynomials are affine linear subspaces, not necessarily vector subspaces.
\item The open-subset dimension lemma and the quasi-affine dimension corollary were rewritten.
\item Homogenisation/dehomogenisation was corrected by fixing a degree; the false global bijection between all affine polynomials and all homogeneous polynomials was removed.
\item The Gr\"obner basis section was rewritten: Dickson's lemma, $S$-polynomials, the division algorithm, and the homogenised Gr\"obner basis theorem now have the correct statements.
\item The Segre embedding proof was corrected to avoid division by possibly zero coordinates, using the identity $P_{ij}P_{ab}=P_{ib}P_{aj}$.
\item The regular-function discussion was corrected: regular functions on irreducible varieties determine rational functions, although not necessarily by one global quotient expression on every chart.
\item The affine morphism adjunction now has the correct target $\mathcal O(X)$.
\item Affine and projective products are separated, and the product variety is distinguished from the product topological space.
\item Rational maps are required to be defined on nonempty open subsets, and the function-field classification is stated as an anti-equivalence.
\item The hypersurface birationality proof now chooses $x_1,\ldots,x_r,y\in K(X)$, not elements of $k$.
\item The projective Jacobian criterion now says rank $n-r$ gives nonsingularity, and the proof uses Euler's relation rather than a false zero-column argument.
\item The singular locus proof was replaced with the standard Jacobian-minor closedness argument and generic smoothness for properness.
\item The blowup section now uses the standard definition of $\operatorname{Bl}_0\bb A^n$ and the strict transform $\overline{\pi^{-1}(Y\setminus\{0\})}$, and it mentions both affine charts for plane curves.
\item The Hartshorne solutions section removes the old incorrect solutions and gives corrected working versions of the exercises most affected by the errata.
\end{enumerate}



\begin{thebibliography}{9}
\bibitem{hartshorne} R. Hartshorne, \emph{Algebraic Geometry}, Springer, 1977.
\bibitem{hartshorne_solutions} W. Troiani, \emph{Hartshorne Exercise Solutions}, unpublished notes.
\bibitem{algebra} W. Troiani, \emph{Elementary Commutative Algebra}, unpublished notes.
\bibitem{grobner} D. Cox, J. Little, and D. O'Shea, \emph{Ideals, Varieties, and Algorithms}, Springer.
\bibitem{matsumura} H. Matsumura, \emph{Commutative Ring Theory}, Cambridge University Press.
\end{thebibliography}
\end{document}

